\documentclass[12pt]{article}
\usepackage{ficha}

\begin{document}

\ficha{Juros Simples e Velocidade Média}

\instrucoes{Resolve os exercícios, apresentando os cálculos e as unidades. Em regime de juro simples, o juro é $J = C \times t \times n$, onde $C$ é o capital, $t$ a taxa de juro (por período) e $n$ o número de períodos. A velocidade média é $v = \dfrac{d}{t}$.}

\secao{Calcular o Juro}{Calcula o juro produzido em regime de juro simples.}
\begin{exercicios}
  \item 1\,000 euros, à taxa anual de 3\,\%, durante 1 ano.
  \item 2\,500 euros, à taxa anual de 4\,\%, durante 3 anos.
  \item 800 euros, à taxa anual de 2,5\,\%, durante 2 anos.
  \item 5\,000 euros, à taxa anual de 1,5\,\%, durante 4 anos.
  \item 12\,000 euros, à taxa anual de 3\,\%, durante 6 meses.
  \item 3\,600 euros, à taxa mensal de 0,5\,\%, durante 10 meses.
  \item 1\,500 euros, à taxa anual de 6\,\%, durante 8 meses.
  \item 20\,000 euros, à taxa anual de 2\,\%, durante 5 anos.
\end{exercicios}

\secao{Capital Acumulado e Valores Desconhecidos}{Resolve.}
\begin{exercicios}
  \item Depositaram-se 4\,000 euros à taxa anual de 2\,\% durante 3 anos. Qual é o capital acumulado?
  \item Que capital, à taxa anual de 5\,\%, rende 150 euros de juro num ano?
  \item Que capital, à taxa anual de 4\,\%, rende 360 euros de juro em 3 anos?
  \item Um capital de 2\,000 euros rendeu 120 euros num ano. Qual foi a taxa de juro anual?
  \item Um capital de 6\,000 euros rendeu 540 euros em 3 anos. Qual foi a taxa anual?
  \item Durante quantos anos devem estar aplicados 3\,000 euros, à taxa anual de 2,5\,\%, para render 300 euros?
  \item Ao fim de 2 anos, à taxa anual de 3\,\%, um capital acumulou 5\,300 euros. Qual era o capital inicial?
  \item O Rui pediu emprestados 1\,200 euros à taxa anual de 8\,\%, a pagar ao fim de 18 meses. Quanto vai pagar no total?
\end{exercicios}

\secao{Velocidade Média}{Calcula a velocidade média.}
\begin{exercicios}
  \item Um carro percorre 240\,km em 3 horas.
  \item Um ciclista percorre 45\,km em 1 hora e 30 minutos.
  \item Um atleta corre 100\,m em 10 segundos. Indica o resultado em m/s e em km/h.
  \item Um comboio percorre 315\,km em 2 horas e 15 minutos.
  \item Um peão anda 2\,km em 25 minutos. Indica o resultado em km/h.
  \item Converte 90\,km/h para m/s.
  \item Converte 15\,m/s para km/h.
  \item Qual é mais rápido: um carro a 108\,km/h ou uma chita a 32\,m/s?
\end{exercicios}

\secao{Distância e Tempo}{Resolve.}
\begin{exercicios}
  \item Um carro desloca-se a 80\,km/h durante 2,5 horas. Que distância percorre?
  \item Uma bicicleta a 18\,km/h anda durante 40 minutos. Que distância percorre?
  \item Quanto tempo demora um comboio a 120\,km/h a percorrer 300\,km?
  \item Quanto tempo demora um peão a 5\,km/h a percorrer 3\,km? Indica em minutos.
  \item Um avião voa a 900\,km/h. Quanto tempo demora a fazer 2\,250\,km?
  \item A Ana saiu de casa às 9h15 e chegou ao destino, a 54\,km, às 10h00. Qual foi a velocidade média?
  \item O som propaga-se no ar a cerca de 340\,m/s. Ouviu-se um trovão 5 segundos depois do relâmpago. A que distância caiu?
\end{exercicios}

\secao{Desafios}{Resolve os problemas.}
\begin{exercicios}
  \item O Banco A paga 3\,\% ao ano em juro simples e o Banco B paga 0,26\,\% ao mês. Qual rende mais para 10\,000 euros durante um ano? Quanto mais?
  \item Um capital duplica ao fim de quantos anos, à taxa anual de 5\,\% em juro simples?
  \item Um carro fez 120\,km a 60\,km/h e depois 120\,km a 120\,km/h. Qual foi a velocidade média na viagem completa? (Atenção: não é 90\,km/h!)
  \item Dois carros partem ao mesmo tempo de cidades a 350\,km de distância, um em direção ao outro, a 80\,km/h e 60\,km/h. Ao fim de quanto tempo se cruzam?
  \item Um ciclista vai de A para B a 20\,km/h e regressa a 30\,km/h. A viagem total demorou 5 horas. Qual é a distância de A a B?
  \item Um empréstimo de 4\,000 euros foi pago ao fim de 2 anos com 4\,480 euros. Qual era a taxa de juro anual?
\end{exercicios}

\end{document}
